

%*****************************************
\chapter{Evaluation}
\label{ch:evaluation}
%*****************************************
\hint{This chapter should describe how the evaluation of the implemented mechanism was done. \\ \\
1. Which evaluation method is used and why? Simulations, prototype? \\
2. What is the goal of the evaluation? Comparison? Proof of concept? \\
3. Wich metrics are used for characterizing the performance, costs, fairness, and efficiency of the system?\\
4. What are the parameter settings used in the evaluation and why? If possible always justify why a certain threshold has been chosen for a particular parameter.  \\
5. What is the outcome of the evaluation? \\
6. It is good scientific practice to separate Methodology from Results and from Discussion. Try to be as objective as possible when describing your results. Comparison to literature, own interpretations and speculations should be left to the Discussion (in this example this could be Section~\ref{sec:analysis} or  \ref{ch:closure}).}
\section{Goal and Methodology}

\section{Evaluation Setup}

\begin{table}[htb]
	\centering
	\caption{Example table. Table captions are usually at the top. Provides examples for the use of multi-columns and multi-rows, using the booktab package style.}
	%%
	\begin{tabular}{rlcr}
		\toprule
		& \multicolumn{3}{c}{Values} \\
		\cmidrule{2-4} % make some line over column 2 and 3 only
		& \textbf{Column left-aligned} & \textbf{Column center-aligned} & \textbf{Column right-aligned} \\		
		\midrule
		
		\multirow{3}{*}{Parameters} 	& P1 & V1 & U1 \\		
										& P2 & V2 & U2 \\	
										& P3 & V3 & U3 \\	
		%%
		\midrule		
		Others	& P4 & \multicolumn{2}{c}{centered multi-column} \\
		\bottomrule		
	\end{tabular}	
	\label{tab:eval_params}
\end{table}

\section{Evaluation Results}

\section{Analysis of Results}
\label{sec:analysis}

